\documentclass{texlr}

\title{Search indexing migration}
\author{Implementation agent}
\date{30 August 2026}
\project{Atlas Platform}
\status{Ready for review}
\repository{https://github.com/example/atlas}
\revision{7f4c2de}

\begin{document}
\maketitle

\section{Executive summary}

The indexing pipeline now writes through a versioned adapter, backfills in
bounded batches, and exposes a reversible cutover. Existing callers retain the
same contract.

\begin{callout}[Reviewer focus]
Confirm that production credentials grant write access to the new index before
changing the read alias.
\end{callout}

\section{Architecture}

\begin{figure}[ht]
  \centering
  \graphviz[width=0.82\linewidth]{diagrams/architecture.dot}
  \caption{Data flow after the migration.}
  \label{fig:architecture}
\end{figure}

The application boundary in Figure~\ref{fig:architecture} keeps the search
provider replaceable and makes rollback independent of the backfill.

\section{Cutover sequence}

\begin{figure}[ht]
  \centering
  \mermaid[width=0.95\linewidth]{diagrams/cutover.mmd}
  \caption{Operational cutover and verification.}
\end{figure}

\decision{Keep the old index for seven days}{The retention window covers the
highest-volume weekly job and permits a low-risk alias rollback.}

\section{Verification}

\begin{enumerate}
  \item Run the consistency check against both indexes.
  \item Compare document counts and sampled result ordering.
  \item Move the read alias and observe error rate for 30 minutes.
\end{enumerate}

\begin{lstlisting}[language=bash,caption={Agent verification command}]
./bin/search verify --old products-v1 --new products-v2
\end{lstlisting}

\end{document}
